% concatenated from main.tex

\documentclass[12pt]{article}

% Packages
\usepackage{lipsum} % For generating dummy text, remove in the final version
\usepackage{hyperref} % For hyperlinks
\usepackage{authblk} % For author affiliations
\usepackage{geometry} % For setting page margins
\usepackage{graphicx}
\usepackage{multirow}
\usepackage{amsmath,amssymb,amsfonts}
\usepackage{amsthm}
\usepackage{mathrsfs}
\usepackage[title]{appendix}
\usepackage{xcolor}
\usepackage{textcomp}
\usepackage{manyfoot}
\usepackage{booktabs}
\usepackage{algorithm}
\usepackage{algorithmicx}
\usepackage{algpseudocode}
\usepackage{listings}
\usepackage[numbers,sort&compress]{natbib}

% Page geometry
\geometry{margin=1in}

% Custom command for keywords
\newcommand{\keywords}[1]{\par\noindent\textbf{Keywords:} #1}

% Theorem environments
\theoremstyle{plain}
\newtheorem{theorem}{Theorem}[section]
\newtheorem{proposition}[theorem]{Proposition}
\newtheorem{note}[theorem]{Note}
\newtheorem{corollary}[theorem]{Corollary}
\newtheorem{lemma}[theorem]{Lemma}
\newtheorem{assumption}[theorem]{Assumption}
\newtheorem{solution}[theorem]{Solution}
\theoremstyle{definition}
\newtheorem{example}{Example}[section]
\newtheorem{remark}{Remark}[section]
\newtheorem{definition}{Definition}[section]

\raggedbottom

% Title and Author
\title{Memory-Type Null Controllability of Heat Equations with Delay Effects}

\author[1]{Dev Prakash Jha\thanks{Corresponding author: \texttt{devprakash.22@res.iist.ac.in}}}
\author[1]{Raju K. George\thanks{\texttt{george@iist.ac.in}}}

\affil[1]{Department of Mathematics, Indian Institute of Space Science and Technology, Thiruvananthapuram, 695547, Kerala, India}

\date{\today}

\begin{document}

\maketitle

\begin{abstract}
This article is devoted to the study of null controllability for evolution equations that incorporate both memory and delay effects. The problem is particularly challenging due to the presence of memory integrals and delayed states, which necessitate strengthening the classical controllability requirement to ensure complete rest at the final time. To address this, we adopt the notion of Delay and memory-type null controllability, which demands the vanishing of the state, the accumulated memory term, and the influence of delay at the terminal time. Utilizing duality arguments, we reduce the controllability analysis to proving suitable observability inequalities for the corresponding adjoint system. We begin with finite-dimensional systems and establish rank-type conditions characterizing controllability. These insights are then extended to parabolic partial differential equations with delay and memory terms. By leveraging Carleman estimates and time-dependent control strategies, we derive sufficient conditions under which controllability can be achieved. Numerical simulations validate the theoretical results and illustrate the critical role of moving control regions in neutralizing the effects of memory and delay.
\end{abstract}



\keywords{Evolution equation with memory, Delay and memory-type null controllability, rank condition, Carleman estimates, observability estimate}





\section{Introduction}
\label{sec:intro}

The study of controllability in evolution equations has a long-standing history, beginning with finite-dimensional linear systems (see \cite{kalman1960general}). In these systems, controllability is typically characterized using algebraic conditions based on the rank of matrices that define the system dynamics and control mechanisms. Over the years, this theory has been extended to a broader class of systems, including those in infinite-dimensional spaces, as well as to nonlinear and stochastic models (see, e.g., \cite{avdonin1995families,fursikov1996controllability,lions1988controlabilite,russell1978controllability,zuazua2007controllability}, and references therein).

Nevertheless, many existing studies address evolution equations that incorporate memory effects, which are often motivated by physical considerations. For instance, in \cite{gurtin1968general}, a modification of Fourier’s law was introduced to resolve the nonphysical implication of instantaneous heat propagation found in the classical heat equation (cf. \cite{cattaneo1958form}), resulting in a heat equation with memory effects.





\begin{equation}\label{eq:Intro_eq_1}
\left\{
\begin{aligned}
   y_t - \sum_{i,j=1}^{n} \left\{ a^{ij}(x) \left[ a y_{x_i} + \int_0^t b(t-s,x)y_{x_i}(s,x)\,ds \right] \right\}_{x_j} = u\chi_{\omega}(x) \quad \text{in } Q, \\
    y = 0 \quad \text{on } \Sigma, \\
    y(0) = y_0 \quad \text{in } \Omega.
\end{aligned}
\right.
\end{equation}

In the system above, \( b(\cdot,\cdot) \) denotes a smooth memory kernel, and \( a \in \{0,1\} \) is a fixed parameter. The domain \( \Omega \subset \mathbb{R}^n \), with \( n \in \mathbb{N} \), is a bounded open set with a boundary \( \partial \Omega \) that is infinitely differentiable. The spatial variable is denoted by \( x = (x_1, \ldots, x_n) \), and \( (a^{ij}(x))_{n \times n} \) is a uniformly positive definite matrix with adequate regularity. The control is localized within a nonempty open subset \( \omega \subset \Omega \). We define \( Q = (0,T) \times \Omega \) and \( \Sigma = (0,T) \times \partial \Omega \), where \( T > 0 \) is a finite time horizon. The symbol \( \chi_{\omega} \) stands for the characteristic function of \( \omega \), and \( \nu = \nu(x) \) is the outward unit normal vector at \( x \in \partial \Omega \).

The existence, uniqueness, and propagation properties of such models have been explored in \cite{yong2011heat}.

When the memory effect is omitted, i.e., \( b(\cdot) \equiv 0 \), and \( a = 1 \), the system simplifies to the standard heat equation:
\begin{equation}\label{eq:Intro_eq_2}
\left\{
\begin{aligned}
    y_t - \sum_{i,j=1}^{n} \left[ a^{ij}(x)y_{x_i} \right]_{x_j} &= u\chi_{\omega}(x) \quad \text{in } Q, \\
    y &= 0 \quad \text{on } \Sigma, \\
    y(0) &= y_0 \quad \text{in } \Omega.
\end{aligned}
\right.
\end{equation}

The null controllability of this classical formulation has been thoroughly studied. For example, it is established in \cite{fursikov1996controllability} that for any \( T > 0 \) and any nonempty open set \( \omega \subset \Omega \), equation (\ref{eq:Intro_eq_2}) is null controllable in \( L^2(\Omega) \). That is, for any initial state \( y_0 \in L^2(\Omega) \), there exists a control input \( u \in L^2((0,T) \times \omega) \) such that the corresponding weak solution \( y(\cdot) \in C([0,T]; L^2(\Omega)) \cap C((0,T]; H_0^1(\Omega)) \) satisfies
\begin{equation}\label{eq:Intro_eq_3}
    y(T) = 0.
\end{equation}

In the context of parabolic equations, it is well-known that due to the infinite speed of propagation, the controllability time $T$ and the control region $\omega$ can be chosen arbitrarily small.

However, this favorable null controllability property does not extend straightforwardly to models with memory effects, where the dynamics are more intricate and less understood (see, e.g., \cite{guerrero2013remarks,halanay2012lack}).

When $a = 0$, it has been demonstrated under certain conditions in \cite{yong2011heat} that system (\ref{eq:Intro_eq_1}) exhibits finite speed of propagation for finite heat pulses. This behavior aligns more closely with physical heat conduction and has meaningful implications in control theory. Specifically, assuming suitable conditions on the coefficients $a^{ij}(\cdot)$ and on the geometry of the control region $\omega$, and provided that $T > 0$ is sufficiently large, the system is controllable in the sense that for any initial and target states $y_0, y_1 \in L^2(\Omega)$, there exists a control function $u \in L^2((0,T) \times \omega)$ such that the corresponding solution $y \in C([0,T]; L^2(\Omega))$ satisfies $y(T) = y_1$ in $\Omega$ \cite{fu2009controllability}. Related discussions can be found in \cite{castro2013exact,russell1978controllability,liu1999rapid,zhang2010unified}.

It is important to note, however, that this type of result does not address the state of the accumulated memory at time $t = T$. Hence, even if $y(T) \equiv 0$, the system cannot be guaranteed to reach full rest unless the memory term also vanishes:
\begin{equation}\label{eq:Intro_eq_4}
    \int_0^T b(T - s, x) y_{x_i}(s,x)\, ds \equiv 0.
\end{equation}

Thus, the controllability described above should be interpreted as partial controllability, since complete control would require steering both the state and the memory term to zero.

In the case where $a = 1$, equation (\ref{eq:Intro_eq_1}) corresponds to a controlled heat equation with a parabolic memory kernel. Recent studies (see \cite{guerrero2013remarks,halanay2012lack,pandolfi2013boundary,zhou2014interior}) have shown that null controllability may fail when the memory kernel $b(\cdot,\cdot)$ is a nontrivial constant and the control region $\omega$ remains fixed over time. Nonetheless, approximate controllability remains achievable under specific settings (see \cite{zhou2014interior,barbu2000controllability}). The complete characterization of this behavior remains an open question, which this work seeks to address.

The main contributions of this paper are twofold: first, to introduce an appropriate definition of controllability for systems involving memory; and second, to demonstrate that, although null controllability may not hold with static control supports, it can be achieved when the support of the control moves in time and eventually covers the entire spatial domain. This approach aligns with the strategy used in \cite{chaves2014null} for viscoelastic systems.

To explore this connection between viscoelastic models and the memory-dependent systems under consideration, we begin by examining the simplest illustrative case:

\begin{equation}\label{eq:Intro_eq_5}
\left\{
\begin{aligned}
    y_t - \Delta y + \int_0^t y(s)\,ds &= u\chi_\omega(x) && \text{in } Q, \\
    y &= 0 && \text{on } \Sigma, \\
    y(0) &= y_0 && \text{in } \Omega.
\end{aligned}
\right.
\end{equation}

Introducing the auxiliary variable $z(t) = \int_0^t y(s)\,ds$, the system can be reformulated as
\begin{equation}
\label{eq:Intro_eq_6}
\left\{
\begin{aligned}
    y_t - \Delta y + z &= u\chi_\omega(x) && \text{in } Q, \\
    z_t &= y && \text{in } Q, \\
    y = z &= 0 && \text{on } \Sigma, \\
    y(0) = y_0,\, z(0) &= 0 && \text{in } \Omega.
\end{aligned}
\right.
\end{equation}

This reformulated system couples a parabolic partial differential equation with an ordinary differential equation, a structure commonly encountered in viscoelastic models (refer to \cite{chaves2014null}). Achieving full null controllability entails steering both the state $y$ and the memory term $z = \int_0^t y(s)\,ds$ to zero at the final time $t = T$. However, due to the influence of the ODE component, controllability cannot be achieved if the control is restricted to a fixed subregion $\omega$ of $\Omega$. Consequently, the control support must traverse the domain throughout the control interval to ensure system controllability.

As will be demonstrated, the strategies and analytical methods outlined in \cite{chaves2014null,chaves2017controllability} can be effectively extended to address this modified framework.

















































To formulate our main results, we consider the following abstract framework:
\begin{equation}
\label{eq:Intro_eq_7}
\left\{
\begin{aligned}
    y_t(t) &= Ay(t) + A_1 y(t-h) + \int_0^t M(t - s)y(s)\,ds + B(t)u(t), \quad t \in [0, T], \\
    y_0(\theta) &= \varphi(\theta),\quad \theta \in [-h, 0].
\end{aligned}
\right.
\end{equation}
Here, $y(\cdot): [-h, \tau] \to Y$ represents the state variable, and $u(\cdot): [0, \tau] \to U$ denotes the control function, where $Y$ and $U$ are Hilbert spaces. Let $C$ denote the Banach space of all continuous functions from $[-h, 0]$ into $Y$, equipped with the supremum norm. The operator $B: U \to Y$ is bounded and linear. The operator $A: D(A) \subset Y \to Y$ is closed with a dense domain $D(A)$ in $Y$, and it generates a $C_0$-semigroup $\{T(t)\}_{t \geq 0}$ on $Y$. Furthermore, $A_1: Y \to Y$ is a bounded linear operator, $\varphi \in C$, $M(\cdot) \in L^1(0,T; \mathcal{L}(Y))$, and $B(\cdot) \in L^2(0,T; \mathcal{L}(U;Y))$.

It can be verified (cf. \cite{yong2011heat}) that, under suitable regularity assumptions on the coefficients $a^{ij}(\cdot)$ and $b(\cdot,\cdot)$, equation \eqref{eq:Intro_eq_1} is a particular case of the general form \eqref{eq:Intro_eq_7}.

Similar to systems without Delay and Memory components, system~\eqref{eq:Intro_eq_7} is said to be \emph{null controllable} if, for every initial state $y_0 \in Y$, there exists a control $u(\cdot) \in L^2(0,T;U)$ such that the corresponding solution $y(\cdot)$ satisfies $y(T) = 0$.

Chaves-Silva~\cite{chaves2017controllability} introduced the notion of \emph{Delay and memory-type null controllability} in the case where $A_1 = 0$, noting that the presence of the memory term $\int_0^t M(t-s)y(s)\,ds$ renders the condition $y(T) = 0$ alone insufficient.

To ensure the system reaches the null state, it is necessary to impose the additional condition:
\begin{equation}
\label{eq:Intro_eq_8}
\int_0^T M(T - s)y(s)\,ds = 0.
\end{equation}
To ensure that the system not only reaches the rest state at time $T$, but also remains at rest for all $t \geq T$, we must impose conditions that eliminate the effects of both the memory and the delay after time $T$. These requirements are:

\begin{equation}
\label{eq:memory_delay_condition}
  y(T) = 0 , \quad \int_{\theta}^{T-h} \widetilde{M}(T - s)y(s)\,ds = 0 \quad \text{for all } \theta \in [-h, 0] \quad \text{and}\quad  y(t) = 0 \quad \text{for all } t \in (T - h, T]
\end{equation}

The first two conditions are standard in Delay and memory-type null controllability, ensuring that both the state and the accumulated memory vanish at time $T$. The third condition is necessary due to the presence of the delay term $y(t - h)$: to prevent this term from reactivating the system for $t > T$, the function $y(t)$ must vanish on the interval $[T - h, T]$.


\begin{remark}
    This hybrid condition generalizes the classical notion of null controllability. In the absence of the memory term ($M \equiv 0$), only the delay condition remains. Conversely, if $A_1 = 0$ (no delay), only the memory condition is required.
\end{remark}



To address this, we introduce the concept of Delay and memory-type null controllability.
\begin{definition}\label{def:Def_memory_type of null control}
Let $\widetilde{M}(\cdot) \in L^1(0,T; \mathcal{L}(Y))$ be a given memory kernel, not necessarily identical to $M(\cdot)$ in \eqref{eq:Intro_eq_7}. The system \eqref{eq:Intro_eq_7} is said to be \textbf{Delay and memory-type null controllable} (with respect to $\widetilde{M}(\cdot)$) if for any $y_0 \in Y$, there exists a control $u(\cdot) \in L^2(0,T;U)$ such that the corresponding mild solution $y(\cdot)$ satisfies:
\begin{equation}\label{eq:Intro_eq_9}
\begin{cases}
   (a)\quad y(T) = 0, \\[6pt]
   (b)\quad \displaystyle\int_{\theta}^{T-h} \widetilde{M}(T - s)y(s)\,ds = 0 \quad \text{for all } \theta \in [-h, 0], \\[6pt]
    (c)\quad y(t) = 0 \quad \text{for all } t \in (T - h, T].
\end{cases}
\end{equation}


\end{definition}



The standard notion of (partial) null controllability, as characterized by (\ref{eq:Intro_eq_3}), can be viewed as a special instance of Delay and memory-type null controllability by choosing $\widetilde{M}(\cdot) \equiv 0$. Full controllability, on the other hand, corresponds to the case where $\widetilde{M} = M$.

The central goal of this work is to analyze the Delay and memory-type null controllability of system \eqref{eq:Intro_eq_7}. By applying standard duality techniques, the controllability question is transformed into an observability problem for the corresponding adjoint system (see Proposition \ref{prop:2.1}). While observability estimates are well developed for classical parabolic systems, the presence of memory terms introduces substantial challenges, which have not yet been fully addressed in the literature. Hence, we focus on this class of systems. Even in the linear setting, establishing Delay and memory-type null controllability necessitates the use of controls supported on a proper subregion $\omega \subset \Omega$, implying that $B(\cdot)$ is not surjective (see Remark \ref{rem:2.2} for more details).

\begin{remark}
\label{rem:1.2}
Unlike classical frameworks that explore trajectory controllability—requiring the vanishing of the memory term at the final time—this paper focuses on  Delay and memory-type null controllability, which demands both $y(T) = 0$ and the nullification of the associated memory integral. In linear systems, this weaker notion is sufficient to characterize the full trajectory controllability. Accordingly, our analysis will be centered around the concept of Delay and memory-type null controllability.
\end{remark}



\section{Abstract Duality and Observability Analogue}

In this section, we address the concept of Delay and memory-type null controllability associated with a memory kernel denoted by $\widetilde{M}(\cdot)$.\\

We begin with the consideration of the following system:
\begin{equation}
\label{eq:Dual_1}
\left\{
\begin{aligned}
    y'(t) &= Ay(t) + A_1 y(t - h), && 0 \leq t \leq T, \\
    y_0(\theta) &= \varphi(\theta), && -h \leq \theta \leq 0.
\end{aligned}
\right.
\end{equation}

It is well-established that, for any initial function $\varphi \in C$, the system admits a unique solution $y(t)$. Let us now define an operator $S(t)$ acting on the space $Y$ by:
\[
S(t)\varphi(0) =
\begin{cases}
y(t), & \text{if } t \geq 0, \\
0, & \text{if } t < 0.
\end{cases}
\]

This operator $S(t): Y \to Y$ is referred to as the fundamental solution corresponding to system \eqref{eq:Dual_1}. It satisfies the following integral representation:
\begin{align}\label{eq:Dual_2}
    S(t) &= T(t) + \int_0^t T(t - s) A_1 S(s - h) \, ds, && t > 0, \notag \\
    S(0) &= I, \notag \\
    S(t) &= 0, && -h \leq t < 0.
\end{align}

By applying Gronwall's inequality, one can conclude the existence of a positive constant $M_S$ such that
\[
\|S(t)\| \leq M_S, \quad \text{for all } t \in [0, \tau].
\]

The unique mild solution $y(\cdot) \in C([-h, T]; Y)$ of equation \eqref{eq:Intro_eq_7} can be represented by the integral formulation given in \cite{wang2005approximate}:
\begin{align}\label{eq:Dual_3}
    y(t) &= S(t)\varphi(0) 
    + \int_0^t S(t - s) B u(s) \, ds + \int_0^t S(t - s) \int_0^s M(t - \tau)y(\tau) \, d\tau \, ds, \quad 0 \leq t \leq \tau, \notag \\
    y_0(\theta) &= \varphi(\theta), \quad -h \leq \theta \leq 0.
\end{align}


















We now examine the adjoint system associated with \eqref{eq:Intro_eq_7}, involving the memory kernel $\widetilde{M}(\cdot)$:\footnotemark

\begin{equation}
\label{eq:dual_eq_1}
\left\{
\begin{aligned}
    -w_t(t) &= A^* w(t) + A_1^* w(t + h) + \int_t^T M(s - t)^* w(s)\,ds - \widetilde{M}(T - t)^* z_T, \quad t \in [0, T], \\
    w(T) &= w_T, \\
    w(t) &= 0, \quad t \in (T, T + h].
\end{aligned}
\right.
\end{equation}


where $w_T, z_T \in Y$.

Throughout the following, we let $K$ denote a generic positive constant, which may differ from line to line unless specified otherwise. We now present the main result.


\begin{proposition}
\label{prop:2.1}
The system described by \eqref{eq:Intro_eq_7} is Delay and memory-type null controllable (with the kernel $\widetilde{M}(\cdot)$) if and only if every solution to \eqref{eq:dual_eq_1} satisfies the observability inequality:
\begin{equation}
\label{eq:enhanced_obs}
\|w(\theta)\|_Y^2 + \int_{ - h}^{\theta} \|w(t+h)\|_Y^2 dt+ \int_{T_1 - h}^{T-h} \|w(t+h)\|_Y^2 dt \leq K \int_0^T \|B(t)^* w(t)\|_U^2 dt,
\end{equation}
for all $w_T, z_T \in Y$ and $T_1 \in [T - h, T]$
\end{proposition}

\begin{proof}
We prove both directions.

\textbf{$(\Rightarrow)$ Sufficiency:} Suppose the observability estimate \eqref{eq:enhanced_obs} holds. Define the set
\[
\mathcal{L} := \left\{ B^*(\cdot)w(\cdot) \in L^2(0, T; U) \mid w(\cdot) \text{ solves } \eqref{eq:dual_eq_1} \right\}.
\]
Define the linear functional $\Lambda : \mathcal{L} \to \mathbb{R}$ by
\[
\Lambda(B^*w) := -\langle w(\theta), \varphi(\theta) \rangle_Y - \int_{- h}^{\theta} \langle A_1^*w(t+h), \varphi(t) \rangle_Y dt + \int_{T_1 - h}^{T-h} \langle A_1^* w(t+h), y(t) \rangle_Y dt,
\]
where $\varphi \in C([-h, 0]; Y)$ is the initial history of the state equation. By the observability inequality, $\Lambda$ is bounded on $\mathcal{L}$ and thus extends to all of $L^2(0, T; U)$ by the Hahn--Banach theorem. Then, by Riesz representation, there exists $u \in L^2(0, T; U)$ such that
\begin{equation}\label{eq:eq_Obs_ineq_3}
    \int_0^T \langle B(t)^* w(t), u(t) \rangle_U dt = \Lambda(B^* w), \quad \forall w_T, z_T \in Y.
\end{equation}


Indeed, for any $w_T,z_T \in Y,$ by \eqref{eq:Intro_eq_7} and (\ref{eq:dual_eq_1}), for any $T_1 \in [T-h,T]$ and $\theta \in [-h,0]$ , we have
\begin{align}\label{eq:eq_Obs_ineq_1}
    &\langle y(T_1), w(T_1) \rangle_Y - \langle w(\theta), \varphi(\theta) \rangle_Y \notag \\
&= \int_{\theta}^{T_1} \left\langle \frac{d}{dt} y(t), w(t) \right\rangle_Y \, dt + \int_{\theta}^{T_1} \left\langle y(t), \frac{d}{dt} w(t) \right\rangle_Y \, dt \notag \\
&= \int_{\theta}^{T_1} \left\langle Ay(t) + A_1y(t-h) + \int_{\theta}^{T_1} M(t-s)y(s)\,ds + B(t)u(t), w(t) \right\rangle_Y \, dt \notag \\
&\quad + \int_{\theta}^{T_1} \left\langle y(t), -A^*w(t) - A_1^*w(t+h) - \int_t^T M(s - t)^*w(s)\,ds + \widetilde{M}(T - t)^*z_T \right\rangle_Y \, dt \notag \\
&= \int_{\theta}^{T_1}\left\langle A_1 y(t-h),w(t)\right\rangle_Y dt -\int_{\theta}^{T_1}\left\langle y(t),A_1^*w(t+h)\right\rangle_Y dt \notag\\
& \quad +\left\langle \int_{\theta}^{T_1}\widetilde{M}(T - t)y(t), z_T  \right\rangle_Y dt + \int_{\theta}^{T_1}\left\langle B(t)u(t),w(t)  \right\rangle_U dt \notag\\
&=I_1 +I_2 +I_3+I_4,
\end{align}
where,
\begin{align*}
    &I_1=\int_{\theta}^{T_1}\left\langle A_1 y(t-h),w(t)\right\rangle_Y dt,\\
   & I_2=  -\int_{\theta}^{T_1}\left\langle y(t),A_1^*w(t+h)\right\rangle_Y dt,\\
   &I_3=\int_{\theta}^{T_1}\left\langle y(t), \widetilde{M}(T - t)^*z_T  \right\rangle_Y dt,\\
   &I_4=\int_{\theta}^{T_1}\left\langle B(t)u(t),w(t)  \right\rangle_U dt.
\end{align*}
Now,

 \begin{align*}
     I_1+I_2&= \int_{\theta}^{T_1}\left\langle A_1 y(t-h),w(t)\right\rangle_Y dt -\int_{\theta}^{T_1}\left\langle y(t),A_1^*w(t+h)\right\rangle_Y dt,
 \end{align*}

put $s=t-h$ in $(I_1)$, then
$$I_1+I_2= \int_{\theta-h}^{T_1-h}\left\langle A_1 y(s),w(s+h)\right\rangle_Y ds -\int_{\theta}^{T_1}\left\langle y(t),A_1^*w(t+h)\right\rangle_Y dt,$$
from (\ref{eq:dual_eq_1})
\begin{align}\label{eq:eq_Obs_ineq_2}
    I_1+I_2&= \int_{\theta-h}^{T_1-h}\left\langle  y(t),A_1^*w(t+h)\right\rangle_Y dt -\int_{\theta}^{T_1}\left\langle y(t),A_1^*w(t+h)\right\rangle_Y dt \notag\\
    &=\int_{\theta-h}^{\theta}\left\langle  y(t),A_1^*w(t+h)\right\rangle_Y dt- \int_{T_1-h}^{T_1}\left\langle  y(t),A_1^*w(t+h)\right\rangle_Y dt \notag \\
    &=\int_{\theta-h}^{\theta}\left\langle  y(t),A_1^*w(t+h)\right\rangle_Y dt- \int_{T_1-h}^{T-h}\left\langle  y(t),A_1^*w(t+h)\right\rangle_Y dt
\end{align}

 Combining \eqref{eq:eq_Obs_ineq_1}, \eqref{eq:eq_Obs_ineq_2}, and \eqref{eq:eq_Obs_ineq_3}, we deduce
\[
\langle y(T_1), w(T_1) \rangle_Y - \int_{\theta}^{T_1} \langle \widetilde{M}(T - t)y(t), z_T \rangle_Y dt = 0, \quad \forall w_T, z_T \in Y.
\]
This implies:
\[
y(T_1) = 0, \quad \int_{\theta}^{T-h} \widetilde{M}(T - t)y(t) dt = 0,
\]
for all $T_1 \in [T-h, T]$.

That is,

\[
y(t) = 0, \quad \int_{\theta}^{T-h} \widetilde{M}(T - t)y(t) dt = 0,
\]
for all $t \in [T-h, T]$, as desired.

























\textbf{$(\Leftarrow)$ Necessity:} Suppose the system is Delay and memory-type null controllable. Assume, for contradiction, that the observability inequality \eqref{eq:enhanced_obs} fails. Then, there exists a sequence $(w_T^n, z_T^n) \in Y \times Y$ with corresponding adjoint solutions $w^n(t)$ such that
\[
\|w(\theta)\|_Y^2 + \int_{ - h}^{\theta} \|w(t+h)\|_Y^2 dt+ \int_{T_1 - h}^{T-h} \|w(t+h)\|_Y^2 dt = 1, \quad \int_0^T \|B(t)^* w^n(t)\|_U^2 dt < \frac{1}{n}.
\]
Define the linear functional:
\[
F_n := \langle w^n(\theta), \varphi(\theta) \rangle_Y + \int_{- h}^{\theta} \langle w^n(t+h), \varphi(t) \rangle_Y dt + \int_{T_1- h}^{T-h} \langle w^n(t+h), y(t) \rangle_Y dt.
\]

By null controllability, there exists a control $u \in L^2(0, T; U)$ such that the solution $y(\cdot)$ satisfies  $\int_{\theta}^{T-h} \widetilde{M}(T - s)y(s) ds = 0$, and $y(t) = 0$ on $[T - h, T]$. Using the duality formula again, we find
\[
F_n = -\int_0^T \langle B(t)^* w^n(t), u(t) \rangle_U dt \to 0 \quad \text{as } n \to \infty.
\]

But since $\|w(\theta)\|_Y^2 + \int_{ - h}^{\theta} \|w(t)\|_Y^2 dt+ \int_{T_1 - h}^{T-h} \|w(t)\|_Y^2 dt = 1$, this implies $F_n \not\to 0$ unless $\varphi \quad \textit{and} \quad y \equiv 0$, contradicting the assumption. Hence, the observability estimate \eqref{eq:enhanced_obs} must hold.

\end{proof}


\begin{remark}
\label{rem:2.2}
The concept of Delay and memory-type null controllability can be naturally reinterpreted through the lens of a nonstandard unique continuation property, which is closely related to the observability inequality \eqref{eq:enhanced_obs}. However, realizing such a property becomes notably challenging when the control operator \( B \) does not depend on time—except in the straightforward case where \( B: U \to Y \) is surjective. Notably, the analysis in \cite{guerrero2013remarks} employs spectral techniques to demonstrate that the observability inequality may fail for heat equations with memory effects, particularly when the control support remains static over time.
\end{remark}

This practical limitation underscores the necessity of allowing the control operator \( B \) to vary with time in the context of Delay and memory-type controllability. Such a time-dependent formulation is especially appropriate in the study of partial differential equations (PDEs), where \( B \) is responsible for localizing control influence in both space and time. A typical and effective approach—motivated by models in viscoelasticity (see \cite{chaves2014null})—is to consider moving control regions, which help in satisfying the strong observability condition stated in \eqref{eq:enhanced_obs}.

\begin{remark}
\label{rem:2.3}
The observability inequality in \eqref{eq:enhanced_obs} is crucial not only for characterizing Delay and memory-type null controllability but also for deriving constructive strategies to compute controls, as discussed in \cite{zuazua2007controllability} for PDE systems. Provided that the inequality holds, consider the following quadratic cost functional defined on the terminal data of the adjoint system \eqref{eq:dual_eq_1}:
\end{remark}

\begin{equation}
J(w_T, z_T) = \frac{1}{2} \int_0^T \|B^*(s)v(s)\|_{\mathcal{U}}^2 \, ds + \langle w(0), y_0 \rangle_Y.
\end{equation}

The functional \( J \), defined for terminal states \( (w_T, z_T) \in Y \times Y \), is continuous and convex on that space. Suppose \( J \) attains its minimum at a point \( (w_T^*, z_T^*) \). Then the corresponding control \( u = B^*v \), where \( v \) is the solution to the adjoint system with this optimal terminal data, provides the control input that fulfills the desired condition in \eqref{eq:Intro_eq_9}.

Nevertheless, the existence of a minimizer for \( J \) is not guaranteed a priori. The observability inequality \eqref{eq:enhanced_obs} is rather weak in the sense that it bounds only the initial state \( u(0) \) in \( Y \), without offering control over \( w_T \) or \( z_T \). Therefore, to ensure the well-posedness of the minimization problem, it becomes necessary to consider the completion of \( Y \times Y \) under the norm induced by the quadratic form.

\begin{equation}
\left( \int_0^T \| B^*(t)v(t) \|_{\mathcal{U}}^2 \, dt \right)^{1/2}.
\end{equation}

It is important to note that verifying whether the seminorm above actually defines a norm is a non-trivial matter. This issue has been thoroughly explored in the context of partial differential equations (PDEs); see \cite{zuazua2007controllability} for details. In the case of the wave equation without memory effects, the associated adjoint system \eqref{eq:dual_eq_1} with \( z_T = 0 \) becomes conservative. Consequently, estimating the norm of \( u(0) \) in the space \( Y \) also leads to an observability estimate for \( w_T \). This observability inequality is then instrumental in minimizing the functional \( J \), which, in this case, no longer depends on \( z_T \), over \( w_T \in Y \).

When memory terms are present, the analysis becomes more delicate. An estimate for \( u(0) \) in \( Y \) does not necessarily provide a bound for \( u \) in the same space. Nevertheless, the backward uniqueness properties of parabolic equations enable us to identify suitable weighted spaces in which the functional \( J \) becomes coercive. This is crucial for ensuring the existence of a minimizer with respect to \( w_T \) within that space.

In the framework of null controllability for systems with memory, such an abstract observability-based formulation of the controllability condition is essential. It serves as a foundation for variational approaches to control construction via functional minimization.

Moreover, this type of characterization finds application in various settings, particularly for developing effective numerical schemes using gradient descent methods.




\section{The Finite-Dimensional Case}
In this section, we study a controlled ordinary differential equation that incorporates a memory term:

\begin{equation}\label{eq:Finite_eq_1}
\left\{
\begin{aligned}
y_t &= Ay(t)+A_1 y(t-h) + \int_0^t M(t-s)y(s)\,ds + Bu(t), \quad t \in (0,T], \\
y(\theta) &= \varphi (\theta), \quad \theta \in [-h, 0].
\end{aligned}
\right.
\end{equation}

Here, \( y = y(t) \) denotes the state vector in \( \mathbb{R}^n \); the matrices \( A, A_1 \in \mathbb{R}^{n \times n} \), and \( M(\cdot) \in L^1(0,T;\mathbb{R}^{n \times n}) \) is the memory kernel. The control input \( u(t) \) takes values in \( \mathbb{R}^m \) for some \( m \in \mathbb{N} \), and \( B \in \mathbb{R}^{n \times m} \) is the control operator. For any matrix \( K \in \mathbb{R}^{n \times m} \), we denote its transpose by \( K^T \).

Building on the formulation introduced in the previous section, we now consider the associated adjoint system corresponding to the memory kernel \( M(\cdot) \in L^1(0,T; \mathbb{R}^{n \times n}) \):

\begin{equation}
\label{eq:Finite_eq_2}
\left\{
\begin{aligned}
    w_t(t) &= -A^* w(t) - A_1^* w(t + h) - \int_t^T M(s - t)^* w(s)\,ds +\widetilde{M}(T - t)^* z_T, \quad t \in [0, T], \\
    w(T) &= w_T, \\
    w(t) &= 0, \quad t \in (T, T + h].
\end{aligned}
\right.
\end{equation}

where \( w_T, x_T \in \mathbb{R}^n \).

We now state the following result:

\begin{theorem}\label{th:Finite_the_1}
    \begin{itemize}
    \item[(i)] Suppose that \( M(\cdot), \widetilde{M}(\cdot) \in L^1(0,T;\mathbb{R}^{n \times n}) \). Then for any solution \( w \) of \eqref{eq:Finite_eq_2}, the condition
    \begin{equation}\label{eq:Finite_eq_3}
         B^T w = 0 \text{ in } [0,T] \quad \Rightarrow \quad x_T = \widetilde{M}(t)^T x_T = 0 \quad \text{for almost every } t \in [0,T]
    \end{equation}
    ensures that the system \eqref{eq:Finite_eq_1} is Delay and memory-type null controllable.

    \item[(ii)] Assume \( M(\cdot) = G M'(\cdot) \) and \( \widetilde{M}(\cdot) = \widetilde{G} M'(\cdot) \) for some constant matrices \( G, \widetilde{G} \in \mathbb{R}^{n \times n} \). If the system \eqref{eq:Finite_eq_1} is Delay and memory-type null controllable, then for every solution of \eqref{eq:Finite_eq_2}, it holds that
     \begin{equation}\label{eq:Finite_eq_4}
         B^T w = 0 \text{ in } [0,T] \quad \Rightarrow \quad \widetilde{M}(t)^T x_T = 0 \quad \text{for almost every } t \in [0,T].
    \end{equation}
    \end{itemize}
\end{theorem}


\begin{proof}
To prove (i), by  using (\ref{eq:Finite_eq_3}), we devide in multipal steps: \\
We aim to prove the inequality:
\begin{equation}\label{eq:obs}
|w_T|^2 + \left( \int_0^T |\widetilde{M}(t)^\top z_T|\, dt \right)^2 \leq C \int_0^T |B^\top w(t)|^2\, dt,
\end{equation}
under the assumption that for any solution $w(t)$ to the adjoint system,
\[
B^\top w(t) \equiv 0 \text{ on } [0, T] \Rightarrow w_T = 0 \text{ and } \widetilde{M}(t)^\top z_T = 0 \text{ a.e. } t \in [0, T].
\]

\textit{Step 1: Contradiction Hypothesis}

Assume, for contradiction, that inequality \eqref{eq:obs} fails. Then for every integer $k \in \mathbb{N}$, there exist terminal data $(w_T^k, z_T^k) \in \mathbb{R}^n \times \mathbb{R}^n$ and corresponding solution $w^k(t)$ of the adjoint system such that:
\[
\int_0^T |B^\top w^k(t)|^2 dt < \frac{1}{k^2} \left( |w_T^k|^2 + \left( \int_0^T |\widetilde{M}(t)^\top z_T^k| dt \right)^2 \right).
\]

\textit{Step 2: Normalize and Pass to the Limit}

Define normalization constants:
\[
N_k := \sqrt{ |w_T^k|^2 + \left( \int_0^T |\widetilde{M}(t)^\top z_T^k| dt \right)^2 } > 0,
\]
and set:
\[
\hat{w}_T^k = \frac{w_T^k}{N_k}, \quad \hat{z}_T^k = \frac{z_T^k}{N_k}.
\]

Let $\hat{w}^k(t)$ denote the solution to the adjoint system with terminal condition $(\hat{w}_T^k, \hat{z}_T^k)$. Then:
\[
\int_0^T |B^\top \hat{w}^k(t)|^2 dt < \frac{1}{k^2}, \quad \text{and} \quad |\hat{w}_T^k|^2 + \left( \int_0^T |\widetilde{M}(t)^\top \hat{z}_T^k| dt \right)^2 = 1.
\]

By compactness in finite dimensions, we may extract a weakly converging subsequence:
\[
\hat{w}_T^k \rightharpoonup w_T^*, \quad \hat{z}_T^k \rightharpoonup z_T^*.
\]

Moreover, the corresponding solutions satisfy:
\[
B^\top \hat{w}^k(t) \to 0 \text{ in } L^2(0,T) \Rightarrow B^\top w^*(t) = 0.
\]

\textit{Step 3: Use the Hypothesis}

By the assumption, $B^\top w^*(t) \equiv 0$ implies:
\[
w_T^* = 0, \quad \widetilde{M}(t)^\top z_T^* = 0 \text{ for all } t.
\]

Thus:
\[
|w_T^*|^2 + \left( \int_0^T |\widetilde{M}(t)^\top z_T^*| dt \right)^2 = 0.
\]

\textit{Step 4: Contradiction}

But by weak lower semi-continuity of the norm:
\[
\liminf_{k \to \infty} \left( |\hat{w}_T^k|^2 + \left( \int_0^T |\widetilde{M}(t)^\top \hat{z}_T^k| dt \right)^2 \right) \geq 1,
\]
contradicting the convergence to zero of the limit. Therefore, the inequality \eqref{eq:obs} must hold.\\


\textit{Step 5: Energy estimation }

We aim to prove the estimate:
\begin{equation}
\label{eq:energy_estimate}
\|w(\theta)\|_Y^2 + \int_{ - h}^{\theta} \|w(t+h)\|_Y^2 dt+ \int_{T_1 - h}^{T-h} \|w(t+h)\|_Y^2 dt \leq C \left[ |w_T|^2_{\mathbb{R}^n} + \left( \int_0^T \widetilde{M}(t)^\top z_T \, dt \right)^2 \right] \quad \forall w_T, z_T \in \mathbb{R}^n.
\end{equation}

We do not attempt a direct energy computation due to the presence of a future-integral memory term in the adjoint system:
\begin{equation}
\label{eq:adjoint}
\dot{w}(t) = -A^\top w(t) - \int_t^T M(s - t)^\top w(s) \, ds + \widetilde{M}(T - t)^\top z_T, \quad w(T) = w_T.
\end{equation}

Instead, we use the well-posedness and stability theory for linear integro-differential equations. From standard results on such systems (see, e.g., Dafermos or Lions), there exists a constant $C > 0$ such that the solution $w(\cdot)$ satisfies the stability estimate:
\[
\|w\|_{C([0, T]; \mathbb{R}^n)} \leq C \left( |w_T|_{\mathbb{R}^n} + \left\| \widetilde{M}(\cdot)^\top z_T \right\|_{L^1(0, T; \mathbb{R}^n)} \right).
\]

In particular, evaluating at $t = 0$ gives:
\[
|w(0)|_{\mathbb{R}^n} \leq C \left( |w_T|_{\mathbb{R}^n} + \int_0^T |\widetilde{M}(t)^\top z_T| \, dt \right).
\]

Applying the Cauchy--Schwarz inequality, we estimate:
\[
\left( \int_0^T |\widetilde{M}(t)^\top z_T| \, dt \right)^2 \leq T \int_0^T |\widetilde{M}(t)^\top z_T|^2 \, dt \leq T \| \widetilde{M} \|_{L^2(0, T)}^2 \cdot |z_T|^2.
\]

Combining the estimates and absorbing constants into $C$, we obtain:
\[
\|w(\theta)\|_Y^2 + \int_{ - h}^{\theta} \|w(t+h)\|_Y^2 dt+ \int_{T_1 - h}^{T-h} \|w(t+h)\|_Y^2 dt \leq C \left( |w_T|^2_{\mathbb{R}^n} + \left( \int_0^T \widetilde{M}(t)^\top z_T \, dt \right)^2 \right),
\]
which completes the proof of inequality \eqref{eq:energy_estimate}.
Combining (\ref{eq:obs}) and (\ref{eq:energy_estimate}), we show that
\begin{equation}\label{eq:Finite_eq_6}
    \|w(\theta)\|_Y^2 + \int_{ - h}^{\theta} \|w(t+h)\|_Y^2 dt+ \int_{T_1 - h}^{T-h} \|w(t+h)\|_Y^2 dt \leq C\int_0^T |B^{T}w(t)|_{\mathbb{R}^m}^2dt \quad \forall w_T,z_T \in \mathbb{R}^n
\end{equation}


Next, we address (ii). From Proposition \ref{prop:2.1} and the null controllability of (\ref{eq:Finite_eq_1}), we conclude that any solution of (\ref{eq:Finite_eq_2}) satisfies the energy bound in (\ref{eq:Finite_eq_6}). Thus, \( w(\theta) = 0 \). Since \( B^T w \equiv 0 \) on \( [0,T] \), define \( \varphi = w_t \). Referring to the first equation in (\ref{eq:Finite_eq_2}) and using \( M(\cdot) = G M'(\cdot) \), \( \widetilde{M}(\cdot) = \widetilde{G} M'(\cdot) \), we obtain:
\begin{align*}
\varphi_t &= -A^T \varphi(t)-A_1^T \varphi(t+h) + M(0)^T w + \int_t^T M(s - t)^T w(s)ds - \widetilde{M}(T - t)^T x_T \\
&= -A^T \varphi (t)-A_1^T \varphi(t+h) + M(0)^T w + \int_t^T M(s - t)^T w(s)ds - \widetilde{M}(T - t)^T x_T \\
&= -A^T \varphi(t)-A_1^T \varphi(t+h) - \int_t^T M(s - t)^T \varphi(s)ds + \widetilde{M}(T - t)^T (G^T w_T - \widetilde{G}^T x_T).
\end{align*}

Thus, \( \varphi \) satisfies:
\begin{equation}
\label{eq:Finite_eq_8}
\left\{
\begin{aligned}
\varphi_t &= -A^T \varphi(t)-A_1^T \varphi(t-h) - \int_t^T M(s - t)^T \varphi(s)ds + \widetilde{M}(T - t)^T (G^T w_T - \widetilde{G}^T x_T), \quad t \in [0, T), \\
\varphi(T) &= -A^T w_T-A^T w(T-h) + \widetilde{M}(0)^T z_T.
\end{aligned}
\right.
\end{equation}

Recognizing that (\ref{eq:Finite_eq_8}) shares the same structure as (\ref{eq:Finite_eq_2}), the energy estimate (\ref{eq:Finite_eq_6}) implies
\begin{equation}
\label{eq:Finite_eq_9}
 \|\varphi(\theta)\|_Y^2 + \int_{ - h}^{\theta} \|\varphi(t+h)\|_Y^2 dt+ \int_{T_1 - h}^{T-h} \|\varphi(t+h)\|_Y^2 dt \leq C \int_0^T |B^T \varphi(s)|^2 ds = C \int_0^T |B^T w_t(s)|^2 ds = 0.
\end{equation}

Consequently, \( w_t(\theta) = 0 \). Iterating this argument shows that all time derivatives vanish at \( t = \theta \), i.e., \( \frac{d^k w(t)}{dt^k}\big|_{t=\theta} = 0 \) for all \( k \in \mathbb{N}_0 \). Since \( w(t) \) is analytic in time, this leads to \( w(t) \equiv 0 \) on \( [0,T] \). Hence,
\[
w_T = \widetilde{M}(t)^T z_T = 0 \quad \text{for all } t \in [0, T].
\]
\end{proof}


As a direct implication of Theorem \ref{th:Finite_the_1}, we obtain the following corollary:
\begin{corollary}
    \textit{
Suppose that \( M(\cdot) = M \in \mathbb{R}^{n \times n} \), \( \widetilde{M}(\cdot) = \widetilde{M} \in \mathbb{R}^{n \times n} \), and there exists \( G, \widetilde{G} \in \mathbb{R}^{n \times n} \) such that \( M = G M' \) and \( \widetilde{M} = \widetilde{G} M' \). Then, system (\ref{eq:Finite_eq_1}) is Delay and memory-type null controllable with kernel \( M \) if and only if every solution to (\ref{eq:Finite_eq_2}) satisfying
\[
B^T w \equiv 0 \quad \text{on } [0,T]
\]
also satisfies
\[
w_T = \widetilde{M}^T z_T = 0.
\]
}
\end{corollary}

We now derive certain rank conditions that characterize the Delay and memory-type null controllability of system (\ref{eq:Finite_eq_1}).

\begin{remark}

The Delay and memory-type null controllability can also be addressed for the more general system:
\begin{equation}
\label{eq:Finite_eq_34}
\left\{
\begin{aligned}
& y_t = A(t)y + \int_0^t M(t-s)y(s)\,ds + B(t)u, && t \in (0,T], \\
& y(0) = y_0,
\end{aligned}
\right.
\end{equation}
where \( A \in L^\infty(0,T;\mathbb{R}^{n \times n}) \) and \( B \in L^\infty(0,T;\mathbb{R}^{n \times m}) \). For this case, one can analyze the extended system:
\begin{equation}
\label{eq:Finite_eq_35}
\left\{
\begin{aligned}
& y_t = A(t)y + \int_0^t M(t-s)y(s)\,ds + B(t)u, && t \in (0,T], \\
& z_t = \widetilde{M}(0)y + \int_0^t \widetilde{M}'(t-s)y(s)\,ds, && t \in (0,T], \\
& y(0) = y_0, \quad z(0) = 0,
\end{aligned}
\right.
\end{equation}
with the terminal condition (\ref{eq:Intro_eq_9}) being equivalent to requiring \( y(T) = z(T) = 0 \).

\vspace{1em}
In the special case where both memory kernels are constant matrices, i.e., \( M(\cdot) \equiv M \in \mathbb{R}^{n \times n} \) and \( \widetilde{M}(\cdot) \equiv \widetilde{M} \in \mathbb{R}^{n \times n} \), and \( M = \widetilde{G}\widetilde{M} \) for some \( G \in \mathbb{R}^{n \times n} \), the extended system \eqref{eq:Finite_eq_35} simplifies to:
\begin{equation}
\label{eq:Finite_eq_36}
\left\{
\begin{aligned}
& y_t = A(t)y + Gz + B(t)u, && t \in (0,T], \\
& z_t = \widetilde{M}y, && t \in (0,T], \\
& y(0) = y_0, \quad z(0) = 0,
\end{aligned}
\right.
\end{equation}
The controllability of this system can be studied using the Silverman–Meadows condition (see, e.g., \cite[Theorem 1.18, p. 11]{coron2007control}). However, when the kernels \( M(\cdot) \) and \( \widetilde{M}(\cdot) \) are not constant, even under the assumption that \( A(\cdot) \equiv A \) is constant, a generalization of the Silverman–Meadows condition to system \eqref{eq:Finite_eq_35} remains unknown. In this context, the result presented in \cite[Theorem 3.3]{chaves2017controllability} can be interpreted as a Delay and memory-type counterpart of the Silverman–Meadows condition, applicable when both \( A(\cdot) = A \) and \( B(\cdot) = B \) are constant matrices.

\end{remark}

\section*{Observability}

We now consider a heat equation with memory and a moving control region $\omega(\cdot) \subset \Omega$, governed by a space-independent memory kernel $M(\cdot) \in L^1(0, T)$:

\begin{equation}
\label{eq:60}
\left\{
\begin{aligned}
    &y_t - \Delta y(t) - \Delta y(t-h) + \int_0^t M(t-s)y(s)\,ds = u\chi_{\omega(t)}(x) && \text{in } Q, \\
    &y = 0 && \text{on } \Sigma, \\
    &y(\theta) = y_0 && -h \leq \theta \leq 0,\quad \text{in } \Omega.
\end{aligned}
\right.
\end{equation}

In the following subsection, we clarify the assumptions imposed on the time-dependent control region and discuss their implications, which are essential for analyzing the controllability of this memory-influenced heat equation.

We assume that the control region $\omega(\cdot)$ is defined via the flow $X(x, t, t_0)$, which describes the evolution of a fixed reference subset under a velocity field 
$f \in C([0, T]; W^{2, \infty}(\mathbb{R}^n; \mathbb{R}^n))$. Specifically, the flow $X$ satisfies:
\[
\left\{
\begin{aligned}
\frac{\partial X(x, t, t_0)}{\partial t} &= f(t, X(x, t, t_0)), \quad t \in [0, T], \\
X(x, t_0, t_0) &= x \in \mathbb{R}^n.
\end{aligned}
\right.
\]

The following structural condition is assumed (see \cite{chaves2014null}):

\begin{assumption}
\label{Ass:Assum_Obser_1}
There exist a flow $X(x, t, 0)$ generated by a vector field $f \in C([0, T]; W^{2, \infty}(\mathbb{R}^n; \mathbb{R}^n))$, a bounded open smooth set $\omega_0 \subset \mathbb{R}^n$, a smooth curve $\Gamma(\cdot) \in C^\infty([0, T]; \mathbb{R}^n)$, and two times $0 \leq t_1 < t_2 \leq T$ such that:
\[
\left\{
\begin{aligned}
&\Gamma(t) \in X(\omega_0, t, 0) \cap \Omega \quad \forall\, t \in [0, T], \\
&\overline{\Omega} \subset \bigcup_{t \in [0, T]} X(\omega_0, t, 0) = \{ X(x, t, 0) \mid x \in \omega_0,\, t \in [0, T] \}, \\
&\Omega \setminus \overline{X(\omega_0, t, 0)} \text{ is nonempty and connected for } t \in [0, t_1] \cup [t_2, T], \\
&\Omega \setminus \overline{X(\omega_0, t, 0)} \text{ has two disjoint nonempty connected components for } t \in (t_1, t_2), \\
&\forall\, \bar{t} \in (t_1, t_2),\, \exists\, \gamma(\cdot) \in C([0, T]; \Omega) \text{ such that } \gamma(\bar{t}) \in X(\omega_0, \bar{t}, 0).
\end{aligned}
\right.
\]
\end{assumption}













To proceed, we utilize a previously established weight function, as presented in the following lemma.

\begin{lemma}\cite{chaves2014null}\label{lemma:lemma_Observation_1}
Let Assumption \ref{Ass:Assum_Obser_1} hold, and let $\omega$ and $\omega_1$ be any two nonempty open subsets of $\mathbb{R}^n$ satisfying $\overline{\omega} \subset \omega_1$ and $\overline{\omega_1} \subset \omega$. Then, there exists a constant $\delta \in (0, T/2)$ and a function $\psi \in C^\infty(\overline{Q})$ such that
\[
\left\{
\begin{array}{ll}
\nabla \psi(t,x) \neq 0, & \text{for } t \in [0,T],~x \in \overline{\Omega} \setminus X(\omega_1,t,0), \\
\psi_t(t,x) \neq 0, & \text{for } t \in [0,T],~x \in \overline{\Omega} \setminus X(\omega_1,t,0), \\
\psi_t(t,x) > 0, & \text{for } t \in [0,\delta],~x \in \overline{\Omega} \setminus X(\omega_1,t,0), \\
\psi_t(t,x) < 0, & \text{for } t \in [T-\delta,T],~x \in \overline{\Omega} \setminus X(\omega_1,t,0), \\
\partial_\nu \psi(t,x) \leq 0, & \text{for } t \in [0,T],~x \in \partial \Omega, \\
\psi(t,x) \geq \frac{3}{4} \|\psi\|_{L^\infty(Q)}, & \text{for } t \in [0,T],~x \in \overline{\Omega}.
\end{array}
\right.
\]
\end{lemma}

Following \cite{chaves2014null}, let us consider a smooth function $g \in C^\infty(0,T)$ defined by
\[
g(t) = 
\begin{cases}
\frac{1}{t} & \text{if } 0 < t < \delta/2, \\
\text{strictly decreasing} & \text{if } 0 < t < \delta, \\
1 & \text{if } \delta \leq t \leq \frac{T}{2}, \\
g(T - t) & \text{if } \frac{T}{2} \leq t < T,
\end{cases}
\]
and define the following weight functions on $Q$:
\[
\varphi(t,x) = g(t) \left( e^{\lambda \|\psi\|_{L^\infty(Q)}} - e^{\lambda \psi(t,x)} \right), \quad
\theta(t,x) = g(t) e^{\lambda \psi(t,x)},
\]
where $\lambda > 0$ is a fixed parameter. For any $p \in H^{1,2}(Q)$, $q \in L^2(Q)$, and $s > 0$, we introduce the quantities:
\begin{align}
I_H(p(t)) &= \int_Q \left[ (s\theta)^{-1}(|\Delta p(t)|^2+|\Delta p(t-h)|^2 + |p_t|^2) + \lambda^2 s \theta |\nabla p|^2 + \lambda^4 (s \theta)^3 |p|^2 \right] e^{-2s\varphi} \, dx \, dt,  \\
I_O(q) &= \lambda^{2s} \int_Q \theta |q|^2 e^{-2s\varphi} \, dx \, dt. 
\end{align}

The subsequent two lemmas are adapted from the results established in \cite{chaves2014null}.
\begin{lemma}\cite{chaves2014null}
\label{lemma:lemma_Observation_2}
Assume that Assumption \ref{Ass:Assum_Obser_1} is satisfied, and let $\omega_1$ be as defined in Lemma \ref{lemma:lemma_Observation_1}. Then, there exist constants $\lambda_0 > 0$ and $s_0 > 0$ such that for any $\lambda \geq \lambda_0$, $s \geq s_0$, and for all functions $p \in C([0,T];L^2(\Omega))$ satisfying $p_t + \Delta p(t) + \Delta p(t-h) \in L^2(0,T;L^2(\Omega))$, the following estimate holds:
\begin{equation}
I_H(p) \leq C \left( \int_Q |p_t + \Delta p(t)+ \Delta p(t-h)|^2 e^{-2s\varphi} \, dx \, dt + \lambda^4 s^3 \int_0^T \int_{X(\omega_1,t,0)} \theta^3 |p|^2 e^{-2s\varphi} \, dx \, dt \right).
\end{equation}
\end{lemma}

\begin{lemma}\cite{chaves2014null}
\label{lemma:lemma_Observation_3}
Suppose Assumption \ref{Ass:Assum_Obser_1} holds and $\omega$ is defined as in Lemma \ref{lemma:lemma_Observation_1}. Then there exist constants $\lambda_1 \geq \lambda_0$ and $s_1 \geq s_0$ such that for all $\lambda \geq \lambda_1$, $s \geq s_1$, and $q \in H^1(0,T;L^2(\Omega))$, the following inequality is satisfied:
\begin{equation}
I_O(q) \leq C \left( \int_Q |q|^2 e^{-2s\varphi} \, dx \, dt + \lambda^2 s^2 \int_0^T \int_{X(\omega,t,0)} \theta^2 |q|^2 e^{-2s\varphi} \, dx \, dt \right).
\end{equation}
\end{lemma}

An immediate consequence of Lemma \ref{lemma:lemma_Observation_3} is the following result.

\begin{corollary}\cite{chaves2017controllability}
\label{cor:cor_1}
Under the assumptions stated in Lemma \ref{lemma:lemma_Observation_3}, for any $\lambda \geq \lambda_1$, $s \geq s_1$, $m \in \mathbb{N}$, and $q \in H^m(0,T;L^2(\Omega))$, the following estimate holds:
\begin{equation}
I_O(q) + \sum_{k=1}^{m-1} I_O(\partial_t^k q)
\leq C \left( \int_Q |\partial_t^m q|^2 e^{-2s\varphi} \, dx \, dt + \int_0^T \int_{X(\omega,t,0)} (\lambda s \theta)^{P(m)} |q|^2 e^{-2s\varphi} \, dx \, dt \right),
\end{equation}
where $P(m)$ denotes a polynomial function of $m$.
\end{corollary}




















































To investigate Delay and memory-type null controllability for equation~\eqref{eq:60} involving the memory kernel $M(\cdot)$, we make use of the dual system derived from~\eqref{eq:dual_eq_1}. Specifically, we consider the following adjoint system associated with~\eqref{eq:60}:
\begin{equation}
\label{eq:74}
\left\{
\begin{aligned}
w_t &= -\Delta w(t) - \Delta w(t-h) - \int_t^T M(s - t)w(s)\,ds + \widetilde{M}(T - t)z_T && \text{in } Q, \\
w &= 0 && \text{on } \Sigma, \\
w(T) &= w_T && \text{in } \Omega,
\end{aligned}
\right.
\end{equation}
where $w_T, z_T \in L^2(\Omega)$ are given terminal data.

In the subsequent analysis, we consider memory kernels of the form
\begin{equation}
\label{eq:75}
M(t) = e^{at}\sum_{k=0}^K a_k t^k, \quad \widetilde{M}(t) = e^{at} \sum_{k=0}^K b_k t^k,
\end{equation}
with $K \in \mathbb{N}$ and coefficients $a_0, \dots, a_K$, $b_0, \dots, b_K$ belonging to $\mathbb{R}$.

We now state an observability result for the solution of the adjoint system~\eqref{eq:74}.

\begin{theorem}\label{The:The_Obs_1}
Suppose Assumption~\ref{Ass:Assum_Obser_1} is satisfied, and let $\omega$ be an open subset of $\Omega$ such that $\overline{\omega} \subset \Omega$. Assume that the memory kernels $M$ and $\widetilde{M}$ are given by~\eqref{eq:75}. Then, the solution $w$ to~\eqref{eq:74} satisfies the following observability estimate:
\begin{equation}
\label{eq:76}
\|w(\theta)\|_{\bar{X}(0)}^2 \leq C \int_0^T \int_\omega |w|^2 \, dx\,dt \quad \forall w_T, z_T \in L^2(\Omega),
\end{equation}
where $w(t) = X(\cdot, t, 0)$ and $C > 0$ is a constant independent of $w_T$ and $z_T$.
\end{theorem}


\begin{proof}
Without loss of generality, we assume $\alpha = 0$ in \eqref{eq:75}. (If $\alpha \neq 0$, one can apply the transformation $w(\cdot) \mapsto e^{-\alpha t} w(\cdot)$ in \eqref{eq:74}.) Define
\[
Z = -\int_t^T M(s - t)w(s)\,ds + \widetilde{M}(T - t)z_T.
\]
Using \eqref{eq:75}, it follows that
\[
\partial_t^{K+1} Z = \sum_{k=0}^K (-1)^k k! a_k \partial_t^{K - k} w.
\]
Substituting this into \eqref{eq:74}, we arrive at the system
\begin{equation}
\label{eq:79}
\begin{cases}
    w_t + \Delta w(t) + \Delta w(t-h) = Z & \text{in } Q, \\
\partial_t^{K+1} Z = \sum_{k=0}^K (-1)^k k! a_k \partial_t^{K - k} w & \text{in } Q, \\
w = 0 & \text{on } \Sigma.
\end{cases}
\end{equation}

We now differentiate the first and third equations in \eqref{eq:79} with respect to time $K+1$ times, and define $v := \partial_t^{K+1} w$. This leads to the system
\begin{equation}
\label{eq:80}
\begin{cases}
    v_t + \Delta v(t) + \Delta v(t-h) = \sum_{k=0}^K (-1)^k k! a_k \partial_t^{K - k} w & \text{in } Q, \\
v = 0 & \text{on } \Sigma, \\
\partial_t^{K+1} w = v & \text{in } Q.
\end{cases}
\end{equation}

Applying Lemma \ref{lemma:lemma_Observation_2} and Corollary \ref{cor:cor_1} to system \eqref{eq:80}, and absorbing lower-order terms, we obtain the estimate
\begin{align}\label{eq:81}
   & I_H(\dot{v}) + I_O(w) + \sum_{i=1}^K I_O(\partial_t^i w) \notag\\
&\leq C \left( \int_0^T \int_{\omega' \cap X(t,0)} \lambda^4 \varphi^3 |\hat{v}|^2 e^{-2\lambda \varphi} \,dx\,dt + \int_0^T \int_{\omega' \cap X(t,0)} (\lambda \varphi)^P |\partial_t^{K+1} w|^2 e^{-2\lambda \varphi} \,dx\,dt \right),
\end{align}
where $\omega'$ is a nonempty open subset of $\mathbb{R}^n$ such that $\overline{\omega} \subset \omega'$, and $P(K+1)$ is a polynomial depending on $K$.

Using the last equation in \eqref{eq:80} and following the same arguments as in the proof of Corollary \ref{cor:cor_1}, we deduce that for any $\varepsilon > 0$,
\begin{align}
\label{eq:82}
& \int_0^T \int_{X(t,0)} \lambda^4 \varphi^3 |\hat{v}|^2 e^{-2\lambda \varphi} \, dx \, dt \notag \\
& \leq \varepsilon \left[ \int_0^T \left( |\varphi|^2 + |\partial_t v|^2 \right) e^{-2\lambda \varphi} \, dx \, dt + I_O(w) + \sum_{i=1}^K I_O(\partial_t^i w) \right] \notag \\
& \quad + \frac{C}{\varepsilon} \int_0^T \int_{\omega' \cap X(t,0)} (\lambda \varphi)^P |\partial_t^{K+1} w|^2 e^{-2\lambda \varphi} \, dx \, dt.
\end{align}

Combining the inequalities \eqref{eq:81} and \eqref{eq:82}, we conclude
\begin{equation}
\label{eq:83}
I_H(\dot{v}) + I_O(w) + \sum_{i=1}^K I_O(\partial_t^i w) \leq C \int_0^T \int_{\omega' \cap X(t,0)} (\lambda \varphi)^P |\partial_t^{K+1} w|^2 e^{-2\lambda \varphi} \, dx \, dt.
\end{equation}

Finally, from \eqref{eq:83} and by employing standard energy estimates for system \eqref{eq:80}, the desired estimate \eqref{eq:76} readily follows. This concludes the proof of Theorem \ref{The:The_Obs_1}.
\end{proof}









By Proposition \ref{prop:2.1}, and as an immediate consequence of Theorem \ref{The:The_Obs_1}, we establish the following Delay and memory-type null controllability result for system (\ref{eq:60}).

\begin{theorem}\label{the:the_obs_2}
Assume the conditions stated in Theorem \ref{The:The_Obs_1} hold. Then, for every initial state $y_0 \in L^2(\Omega)$, there exists a control function $u \in L^2(Q)$ such that the corresponding solution $y(\cdot)$ to system (\ref{eq:60}), where the control is supported as $u(t) = X(\cdot, t, 0)$ for $t \in (0, T)$, fulfills
\[
y(T) = \int_0^T \widetilde{M}(T - s) y(s) \, ds = 0 \quad \text{in } \Omega.
\]
\end{theorem}

\begin{remark}
To derive Theorem \ref{the:the_obs_2}, Assumption \ref{Ass:Assum_Obser_1} plays a crucial role, although it may not represent the most general or minimal set of conditions. For example, the fourth requirement—that the complement of the control region splits into two connected components at times $t \in \{t_1, t_2\}$—is vital for the construction of the fundamental weight function in Lemma \ref{lemma:lemma_Observation_1}. This construction is instrumental in proving null controllability via Carleman estimates with moving control regions. Nonetheless, results such as those in \cite{martin2013null} show that in one-dimensional cases, this condition is not necessary for achieving null controllability of the structurally damped wave equation with moving control. For further insights regarding the sharpness of Assumption \ref{Ass:Assum_Obser_1}, we refer the reader to the discussion in \cite{chaves2014null}.
\end{remark}



\section{Application}
 Consider the heat control system with a delay given by \cite{sukavanam2011approximate} :

 \begin{equation}
\begin{cases}
\frac{\partial y(t,x)}{\partial t} = \frac{\partial^2 y(t,x)}{\partial x^2} + y(t - h,x) + Bu(t,x) + \int_0^t M(t-s)y(s,x) ds, & 0 < t < T, \; 0 < x < \pi, \\
y(t,0) = y(t,\pi) = 0, & 0 \le t \le T, \\
y(t,x) = \varphi(x), & -h \le t \le 0, \; 0 \le x \le \pi.
\end{cases}
\end{equation}


Let $X = L^2(0, \pi)$ and $A = -\frac{d^2}{dx^2}$ with 
\[
D(A) = \left\{ y \in X : y, \; \frac{dy}{dx} \text{ are absolutely continuous}, \; \frac{d^2 y}{dx^2} \in X \text{ and } y(0) = y(\pi) = 0 \right\}.
\]

Let $\psi_n(x) = \left( \frac{2}{\pi} \right)^{1/2} \sin(nx), \; 0 \le x \le \pi, \; n = 1, 2, 3, \dots$. Note that $\psi_n$ is the eigenfunction corresponding to the eigenvalue $\lambda_n = -n^2$ of the operator $A$ and $\{ \psi_n \}$ is an orthonormal base for $X$. Then for $y \in D(A)$ we have \cite{jha2024exact,jha2024existence} 
\[
y = \sum_{n=1}^\infty \langle y, \psi_n \rangle \psi_n, \quad Ay = -\sum_{n=1}^\infty n^2 \langle y, \psi_n \rangle \psi_n.
\]

It is well known that $A$ generates a compact semigroup $T(t)$ \cite{naito1987controllability,jeong1999approximate,jha2025approximate} and that
\[
T(t)y = \sum_{n=1}^\infty e^{-n^2 t} \langle y, \psi_n \rangle \psi_n, \quad n = 1, 2, 3, \dots, \text{ for } y \in X.
\]

Now we define an infinite dimensional control space $U$ by
\[
U = \left\{ u = \sum_{n=2}^\infty u_n \psi_n, \; \sum_{n=2}^\infty u_n^2 < \infty \right\} \quad \text{with norm defined by } \| u \|_U = \left( \sum_{n=2}^\infty u_n^2 \right)^{1/2}.
\]

Then $U$ is a Hilbert space. Define a continuous linear mapping $\tilde{B}$ from $U$ to $X$ by
\[
\tilde{B}u = 2u_2 \psi_1 + \sum_{n=2}^\infty u_n \psi_n, \quad \text{for } u = \sum_{n=2}^\infty u_n \psi_n \in U.
\]

We define bounded linear operator $B : L^2(0,T:U) \to L^2(0,T:X)$ by $(Bu)(t) = \tilde{B}u(t)$.

Then we reformulate the heat control system to the abstract form of a Delay and memory-type semilinear control system with delay (\ref{eq:Intro_eq_7}) in the following way. Put $A_1 = A$, 
\begin{align*}
y' &= Ay(t) + A_1 y(t - h) + Bu(t)  + \int_0^t M(t-s)y(s,x) ds, \quad 0 < t < T \\
y(t) &= \varphi(t), \quad t \in [-h, 0].
\end{align*}


Since $T(t)$ is compact semigroup, then the fundamental solution $S(t)$ is compact.


It can be shown that assumption \ref{Ass:Assum_Obser_1} is valid with similar arguments as in (\ref{eq:Intro_eq_7}).  
Using Theorem \ref{the:the_obs_2}, the control system given by Eq.~(\ref{eq:Intro_eq_7}) is null controllable if assumption (\ref{Ass:Assum_Obser_1}) is satisfied.




\subsection{Numerical Experiment}

We present a numerical example to illustrate the Delay and memory-type null controllability of the delayed heat equation with memory, as established in Theorem~\ref{the:the_obs_2}. Consider the following equation:

\begin{equation}
\label{eq:numerical}
\begin{cases}
\partial_t y(t,x) = \partial_{xx} y(t,x) + y(t - h, x) + \displaystyle\int_0^t e^{-(t - s)} y(s,x)\,ds + u(t,x)\chi_{\omega(t)}(x), & t \in (0,T),\, x \in (0,\pi), \\
y(t,0) = y(t,\pi) = 0, & t \in [0,T], \\
y(t,x) = \sin(x), & t \in [-h,0],\, x \in (0,\pi),
\end{cases}
\end{equation}
where $h = 0.1$ and $T = 1$. The memory kernel is $M(t) = e^{-t}$ and the initial history is given by $\varphi(t,x) = \sin(x)$. The control is applied in a moving spatial region $\omega(t) = [\delta(t), \delta(t)+0.5]$, where $\delta(t) = \frac{(\pi - 0.5)t}{T}$, which ensures that the entire spatial domain is covered over the interval $[0,T]$.

The control function is chosen as
\[
u(t,x) = -t\, e^{-\pi^2 t} \sin(\pi x),
\]
which satisfies the regularity conditions required by the control framework and vanishes at $t=0$, ensuring compatibility with the Delay and memory-type   null controllability goal.

To numerically solve~\eqref{eq:numerical}, we employ the following discretization:
\begin{itemize}
    \item \textbf{Spatial discretization:} Second-order finite difference with $N_x = 50$ points in $[0,\pi]$.
    \item \textbf{Temporal discretization:} Implicit Euler method with $N_t = 200$ steps over $[0,T]$.
    \item \textbf{Delay term:} The value $y(t-h,x)$ is approximated using stored past values at the corresponding delayed time step.
    \item \textbf{Memory term:} The convolution $\int_0^t \widetilde{M}(t-s)y(s,x)\,ds$ is approximated using a discrete sum over past steps with trapezoidal weights.
    \item \textbf{Control implementation:} The control is applied only at grid points $x_j \in \omega(t_n)$ at each time step, based on the analytical formula for $u(t,x)$.
\end{itemize}
\newpage
Figure~\ref{fig:simulation} displays the numerical solution \( y(t,x) \) and its Slice surface at final time \( t = T \), showing that both the solution state and its outer layer approach the \( Z \)-axis for \( t \in [T - h, T + h] \). This observation validates the theoretical result on delay and memory-type null controllability, confirming that all three conditions in Definition~(\ref{def:Def_memory_type of null control}) are satisfied.


\begin{figure}[htbp]
\centering
\begin{minipage}{0.95\textwidth}
  \centering
  \includegraphics[width=\textwidth]{fig_1.jpg}
\end{minipage}
\hfill
\begin{minipage}{0.95\textwidth}
  \centering
  \includegraphics[width=\textwidth]{fig_2.jpg}
\end{minipage}
\caption{The numerical solution \( y(t,x) \) and its Slice surface at final time \( t = T \) of the heat equation with memory and delay are obtained using the control \( u(t,x) = -t\, e^{-\pi^2 t} \sin(\pi x) \) in a moving region \( \omega(t) \).
}
\label{fig:simulation}
\end{figure}
Figure~\ref{fig:simulation_1} illustrates the numerical solution \( y(t,x) \) and its Slice surface at final time \( t = T \). At this moment, only the first two conditions of Definition~(\ref{def:Def_memory_type of null control}) are satisfied. As a result, both the solution and its outer layer do not approach the \( Z \)-axis for \( t \in [T - h, T + h] \) or at \( t = T \). This outcome confirms that all three conditions in Definition~(\ref{def:Def_memory_type of null control}) are necessary to achieve delay and memory-type null controllability.


\begin{figure}[htbp]
\centering
\begin{minipage}{0.95\textwidth}
  \centering
  \includegraphics[width=\textwidth]{fig_3.jpg}
\end{minipage}
\hfill
\begin{minipage}{0.95\textwidth}
  \centering
  \includegraphics[width=\textwidth]{fig_4.jpg}
\end{minipage}
\caption{The numerical solution \( y(t,x) \) and its Slice surface at final time \( t = T \) of the heat equation with memory and delay are obtained using the control \( u(t,x) = -t\, e^{-\pi^2 t} \sin(\pi x) \) in a moving region \( \omega(t) \).
}
\label{fig:simulation_1}
\end{figure}













































































\section{Data Availability}
	The authors used no available data when preparing it for publication. It does not have any data to share.
	
	
\section{Declarations:}

\textbf{Conflict of interest} \\
All authors have no conflict of interest in any form regarding this paper.















































% References
\bibliographystyle{unsrtnat} % Use numerical references
\bibliography{references} % Replace 'references' with your .bib file

\end{document}




XXXXXXXXXXXXXXXXXXXXXXXXXXXXXXXXXXXXXXXXXXXXXXXXXXXXXXXXXXXXXXXXX

\subsection*{Error Analysis and Convergence}

To demonstrate the accuracy and convergence of the numerical scheme, we perform a grid refinement study. Since the exact solution to~\eqref{eq:numerical} is not known in closed form, we compute a reference solution using a fine spatial discretization ($N_x = 160$) and compare it against coarser grids by interpolating the fine solution.

The error is measured using the $L^2$-norm at the final time $t = T$:
\[
E_{L^2} = \left( \sum_j |y^{\text{coarse}}_j(T) - y^{\text{fine}}(x_j,T)|^2 \Delta x \right)^{1/2}.
\]

The results of the convergence study are summarized in Table~\ref{tab:convergence}, showing a clear decrease in error with mesh refinement, and an observed convergence rate close to 2, consistent with second-order accuracy in space.

\begin{table}[ht]
\centering
\begin{tabular}{cccc}
\toprule
$N_x$ & $\Delta x$ & $L^2$ Error & Observed Rate \\
\midrule
20 & 0.1571 & $1.53 \times 10^{-2}$ & -- \\
40 & 0.0785 & $3.87 \times 10^{-3}$ & 1.98 \\
80 & 0.0393 & $9.75 \times 10^{-4}$ & 1.99 \\
\bottomrule
\end{tabular}
\caption{Error and convergence rate of the numerical scheme at final time $T=1$.}
\label{tab:convergence}
\end{table}




XXXXXXXXXXXXXXXXXXXXXXXXXXXXXXXXXXXXXXXXXXXXXXXXXXXXXXXXXXXXXXXXXX

\begin{theorem}
   \begin{itemize}
    \item[(i)] \textit{Assume that \( M(\cdot), \widetilde{M}(\cdot) \in L^1(0, T; \mathbb{R}^{n \times n}) \cap C^\infty([0, T]; \mathbb{R}^{n \times n}) \), and define sequences \( A_i \), \( M_i(\cdot) \), and \( \widetilde{M}_i(\cdot) \) (for \( i = 1, 2, \ldots \)) inductively by:}
    
    \begin{equation}
    \label{eq:Finite_eq_10}
    \left\{
    \begin{aligned}
        A_{i+1} &= A A_i + M_i(0), \\
        M_{i+1}(\cdot) &= M(\cdot) A_i + M_i'(\cdot), \\
        \widetilde{M}_{i+1}(\cdot) &= \widetilde{M}(\cdot) A_i + \widetilde{M}_i'(\cdot), \\
        A_1 &= A, \quad M_1(\cdot) = M(\cdot), \quad \widetilde{M}_1(\cdot) = \widetilde{M}(\cdot).
    \end{aligned}
    \right.
    \end{equation}

     \textit{If \footnotemark}
    \begin{equation}
    \label{eq:Finite_eq_11}
    \text{rank} \left(
    \begin{array}{ccccccc}
        B & A_1 B & A_2 B & \cdots & A_i B & A_{i+1} B & \cdots \\
        0 & \widetilde{M}_1(0) B & \widetilde{M}_2(0) B & \cdots & \widetilde{M}_i(0) B & \widetilde{M}_{i+1}(0) B & \cdots
    \end{array}
    \right) = 2n,
    \end{equation}
    
    \textit{then system (\ref{eq:Finite_eq_1}) is memory-type null controllable.}
    \item[(ii)]  \textit{Assume that both \( M(\cdot) \) and \( \widetilde{M}(\cdot) \) are analytic in \( [0,T] \), and define \( A_i \) as in (\ref{eq:Finite_eq_9}) \( (i = 1, 2, \dots) \). Assume that \( \frac{d^i \widetilde{M}(t)}{dt^i}|_{t=0} = \widetilde{M}(0) G_i \) for some \( G_i \in \mathbb{R}^{n \times n} \), and define}
\begin{equation}\label{eq:Finite_eq_12}
F_i = A_i + G_1 A_{i-1} + \cdots + G_{i-1} A_1 + G_i.
\end{equation}
\textit{If}
\begin{equation}\label{eq:Finite_eq_13}
\text{rank} \left(
\begin{array}{ccccccc}
B & A_1 B & A_2 B & \cdots & A_i B & A_{i+1} B & \cdots \\
0 & B & F_1 B & \cdots & F_{i-1} B & F_i B & \cdots
\end{array}
\right) = 2n,
\end{equation}
\textit{then system (\ref{eq:Finite_eq_1}) is memory-type null controllable.}
\item[(iii)]  \textit{Assume that \( M(\cdot) \equiv M \in \mathbb{R}^{n \times n} \) and \( \widetilde{M}(\cdot) \equiv \widetilde{M} \in \mathbb{R}^{n \times n} \), and define \( A_i \) \( (i=1,2,\dots) \) inductively by}
\begin{equation}\label{eq:Finite_eq_14}
\left\{
\begin{aligned}
A_{i+1} &= A A_i + M_i, \\
M_{i+1} &= M A_i, \\
A_1 &= A, \quad M_1 = M.
\end{aligned}
\right.
\end{equation}
\textit{If}
\begin{equation}\label{eq:Finite_eq_15}
\text{rank} \left(
\begin{array}{cccccc}
B & A_1 B & A_2 B & \cdots & A_{2n-1} B & A_{2n} B \\
0 & B & A_1 B & \cdots & A_{2n-2} B & A_{2n-1} B
\end{array}
\right) = 2n,
\end{equation}
\textit{then system (\ref{eq:Finite_eq_1}) is memory-type null controllable. Moreover, if \( \widetilde{M} \neq 0 \), then this condition is also necessary for (\ref{eq:Finite_eq_1}) to be memory-type null controllable.}
\end{itemize}

\footnotetext{Note that equation \ref{eq:Finite_eq_14} is time-independent even if \( M_{k+1}(t) \) depends on \( t \). Because of this, equation \ref{eq:Finite_eq_15} is an algebraic condition.}


 
\end{theorem}


\newpage



\begin{proof}


\textit{(i)} Suppose that for some \( w_T, z_T \in \mathbb{R}^n \), the corresponding solution \( w(\cdot) \) of (\ref{eq:Finite_eq_2}) satisfies \( B^\top w(\cdot) \equiv 0 \) in \( [0, T] \).\\
It is easy to see that \( B^\top w(\cdot) \equiv 0 \) in \( [0, T] \) implies:
\begin{equation}\label{eq:Finite_eq_16}
B^\top w_T = 0,
\end{equation}
and
\begin{equation}\label{eq:Finite_eq_17}
0 = -B^\top w_t = B^\top \left( A_1^\top w + \int_t^T M_1(s - t) w(s) ds - \widetilde{M}_1(T - t)^\top z_T \right), \quad \text{in } [0, T].
\end{equation}
By this \ref{eq:Finite_eq_17} and using (\ref{eq:Finite_eq_2}), we find that
\begin{equation}\label{eq:Finite_eq_18}
B^\top A_1^\top w_T - B^\top \widetilde{M}_1(0)^\top z_T = 0
\end{equation}
and
\begin{align}
w_{tt} &= -A_1^\top w_t + M_1(0)^\top w_T + \int_t^T M_1'(s - t)^\top w(s) ds - \widetilde{M}_1'(T - t)^\top z_T \notag \\
&= A_1^\top A_1^\top w + \int_t^T A_1^\top M(s - t)^\top w(s) ds - A_1^\top \widetilde{M}(T - t)^\top z_T \notag \\
&\quad + M_1(0)^\top w_T + \int_t^T M_1'(s - t)^\top w(s) ds - \widetilde{M}_1'(T - t)^\top z_T \notag \\
&= A_2^\top w + \int_t^T M_2(s - t)^\top w(s) ds - \widetilde{M}_2(T - t)^\top z_T, \quad \text{in } [0, T].
\end{align}



\textbf{More generally}, we have
\begin{equation}\label{eq:Finite_eq_19}
B^\top A_i^\top w_T - B^\top \widetilde{M}_i(0)^\top z_T = 0
\end{equation}
and
\begin{align}\label{eq:19}
\frac{d^{i+1} w}{dt^{i+1}} 
&= (-1)^{i+1} A_{i+1}^\top w + (-1)^{i+1} \int_t^T M_{i+1}(s - t)^\top w(s) ds \notag \\
&\quad + (-1)^i \widetilde{M}_{i+1}(T - t)^\top z_T, \quad \text{in } [0, T].
\end{align}

By (\ref{eq:Finite_eq_16}), (\ref{eq:Finite_eq_18}), and (\ref{eq:Finite_eq_19}), we end up with
\begin{equation}\label{eq:Finite_eq_20}
(w_T, -z_T^\top) 
\begin{pmatrix}
B & A_1 B & A_2 B & \cdots & A_i B & A_{i+1} B  \cdots \\
0 & \widetilde{M}_1(0) B & \widetilde{M}_2(0) B & \cdots & \widetilde{M}_i(0) B & \widetilde{M}_{i+1}(0) B \cdots
\end{pmatrix}
= 0.
\end{equation}


By (\ref{eq:Finite_eq_11}) and (\ref{eq:Finite_eq_20}), we conclude that \( w_T = z_T = 0 \). Hence, by the first conclusion of Theorem 3.1, we conclude that (\ref{eq:Finite_eq_1}) is memory-type null controllable.

\textbf{(ii)} As in (i), we suppose that for some \( w_T, z_T \in \mathbb{R}^n \), the corresponding solution \( w(\cdot) \) of (\ref{eq:Finite_eq_2}) satisfies \( B^\top w(\cdot) \equiv 0 \) in \( [0, T] \). Then, we have (\ref{eq:Finite_eq_16}). By (\ref{eq:Finite_eq_9}), (\ref{eq:Finite_eq_12}), and \( \frac{d^i \widetilde{M}(t)}{dt^i}|_{t=0} = \widetilde{M}(0) G_i \), it is easy to check that
\begin{equation}\label{eq:Finite_eq_21}
\left\{
\begin{aligned}
\widetilde{M}_1(0) &= \widetilde{M}(0), \\
\widetilde{M}_{i+1}(0) &= \widetilde{M}(0)(A_i + G_1 A_{i-1} + \cdots + G_{i-1} A_1 + G_i) \\
&= \widetilde{M}(0) F_i, \quad i = 1, 2, \dots
\end{aligned}
\right.
\end{equation}

Hence, (\ref{eq:Finite_eq_18}) reads
\begin{equation}\label{eq:Finite_eq_22}
B^\top A_1^\top w_T - B^\top \widetilde{M}(0)^\top z_T = 0,
\end{equation}
and (\ref{eq:Finite_eq_19}) is specialized as
\begin{equation}\label{eq:Finite_eq_23}
B^\top A_i^\top w_T - B^\top F_{i-1}^\top \widetilde{M}(0)^\top z_T = 0, \quad i = 2, 3, \dots
\end{equation}

By (\ref{eq:Finite_eq_16}), (\ref{eq:Finite_eq_22}), and (\ref{eq:Finite_eq_23}), we obtain that
\begin{equation}\label{eq:Finite_eq_24}
(w_T, -z_T^\top \widetilde{M}(0)) 
\begin{pmatrix}
B & A_1 B & A_2 B & \cdots & A_{i+2} B & \cdots \\
0 & B & F_1 B & \cdots & F_i B & F_{i+1} B \cdots
\end{pmatrix}
= 0.
\end{equation}


By (\ref{eq:Finite_eq_13}) and (\ref{eq:Finite_eq_24}), we see that
\begin{equation}\label{eq:Finite_eq_25}
w_T = \widetilde{M}(0) z_T = 0.
\end{equation}

By (\ref{eq:19}), (\ref{eq:Finite_eq_21}), and (\ref{eq:Finite_eq_25}), it follows that
\begin{align}\label{eq:Finite_eq_25'}
\frac{d^{i+1} w}{dt^{i+1}} \bigg|_{t=T} 
&= (-1)^{i+1} A_{i+1}^\top w_T + (-1)^i F_i^\top \widetilde{M}(0)^\top z_T = 0.
\end{align}

Since both \( M(\cdot) \) and \( \widetilde{M}(\cdot) \) are analytic in \( [0, T] \), so is \( w(\cdot) \). By (\ref{eq:Finite_eq_25'}), we conclude that \( w(\cdot) \equiv 0 \) in \( [0, T] \). Hence, by the first equation in (\ref{eq:Finite_eq_2}), \( \widetilde{M}(t)^\top z_T = 0 \) in \( [0, T] \). Now, by the first conclusion of Theorem 3.1, (\ref{eq:Finite_eq_1}) is memory-type null controllable.




(iii) We can use the result in (ii). For the present case, it is easy to check that the $F_i$ defined by (\ref{eq:Finite_eq_12}) is specialized to $F_i = A_i$. We claim that
\begin{equation}\label{eq:Finite_eq_26}
\text{rank} \begin{pmatrix}
B & A_1B & A_2B & \cdots & A_iB & A_{i+1}B & \cdots \\
0 & B & A_1B & \cdots & A_{i-1}B & A_iB & \cdots
\end{pmatrix}
= \text{rank} \begin{pmatrix}
B & A_1B & A_2B & \cdots & A_{2n+1}B \\
0 & B & A_1B & \cdots & A_{2n}B
\end{pmatrix}.
\end{equation}

It is easy to see that (\ref{eq:Finite_eq_14}) is a special case of (\ref{eq:Finite_eq_9}). From (\ref{eq:Finite_eq_14}), we see that
\[
\begin{pmatrix}
A_{i+1} \\
M_{i+1}
\end{pmatrix}
=
\begin{pmatrix}
A & I_n \\
M & 0
\end{pmatrix}
\begin{pmatrix}
A_i \\
M_i
\end{pmatrix}.
\]

Hence,
\begin{equation}\label{eq:Finite_eq_27}
\begin{pmatrix}
A_{i+1} \\
M_{i+1}
\end{pmatrix}
=
\begin{pmatrix}
A & I_n \\
M & 0
\end{pmatrix}^i
\begin{pmatrix}
A \\
M
\end{pmatrix}.
\end{equation}

Denote by $\lambda^{2n} + a_1 \lambda^{2n-1} + a_2 \lambda^{2n-2} + \cdots + a_{2n}$ the characteristic polynomial of 
$\begin{pmatrix} A & I_n \\ M & 0 \end{pmatrix}$, where $a_1, a_2, \ldots, a_{2n} \in \mathbb{R}$. By the Hamilton--Cayley theorem, it follows that
\begin{equation}\label{eq:Finite_eq_28}
\begin{pmatrix}
A & I_n \\
M & 0
\end{pmatrix}^{2n}
+ a_1 \begin{pmatrix}
A & I_n \\
M & 0
\end{pmatrix}^{2n-1}
+ a_2 \begin{pmatrix}
A & I_n \\
M & 0
\end{pmatrix}^{2n-2}
+ \cdots + a_{2n} I_{2n} = 0.
\end{equation}

Combining (\ref{eq:Finite_eq_27}) and (\ref{eq:Finite_eq_28}), we have
\begin{align*}
\begin{pmatrix}
A_{2n+1} \\
M_{2n+1}
\end{pmatrix}
&=
\begin{pmatrix}
A & I_n \\
M & 0
\end{pmatrix}^{2n}
\begin{pmatrix}
A \\
M
\end{pmatrix} \\
&= -a_1 
\begin{pmatrix}
A & I_n \\
M & 0
\end{pmatrix}^{2n-1}
\begin{pmatrix}
A \\
M
\end{pmatrix}
- a_2 
\begin{pmatrix}
A & I_n \\
M & 0
\end{pmatrix}^{2n-2}
\begin{pmatrix}
A \\
M
\end{pmatrix}
- \cdots - a_{2n}
\begin{pmatrix}
A \\
M
\end{pmatrix} \\
&= -a_1 
\begin{pmatrix}
A_{2n} \\
M_{2n}
\end{pmatrix}
- a_2 
\begin{pmatrix}
A_{2n-1} \\
M_{2n-1}
\end{pmatrix}
- \cdots - a_{2n}
\begin{pmatrix}
A_1 \\
M_1
\end{pmatrix}.
\end{align*}

This gives
\begin{equation}\label{eq:Finite_eq_29}
A_{2n+1} = -a_1 A_{2n} - a_2 A_{2n-1} - \cdots - a_{2n} A_1.
\end{equation}

Similarly,
\begin{equation}\label{eq:Finite_eq_30}
A_{2n+2} = -a_1 A_{2n+1} - a_2 A_{2n} - \cdots - a_{2n} A_2.
\end{equation}

Combining (\ref{eq:Finite_eq_29}) and (\ref{eq:Finite_eq_30}), we find that
\begin{equation}\label{eq:Finite_eq_31}
\begin{pmatrix}
A_{2n+2} \\
A_{2n+1}
\end{pmatrix}
= -a_1 
\begin{pmatrix}
A_{2n+1} \\
A_{2n}
\end{pmatrix}
- a_2 
\begin{pmatrix}
A_{2n} \\
A_{2n-1}
\end{pmatrix}
- \cdots - a_{2n}
\begin{pmatrix}
A_2 \\
A_1
\end{pmatrix}.
\end{equation}









Inductively, from (\ref{eq:Finite_eq_31}), one can show that each 
$\begin{pmatrix} A_{k+1} \\ A_k \end{pmatrix}$ ($k \geq 2n + 1$)
can be expressed as a linear combination of 
$\begin{pmatrix} A_2 \\ A_1 \end{pmatrix}, \begin{pmatrix} A_3 \\ A_2 \end{pmatrix}, \ldots, \begin{pmatrix} A_{2n+1} \\ A_{2n} \end{pmatrix}$. Consequently, (\ref{eq:Finite_eq_26}) is verified.

By the result in (ii) and (\ref{eq:Finite_eq_26}), it is easy to see that under the condition (\ref{eq:Finite_eq_15}), (\ref{eq:Finite_eq_1}) is memory-type null controllable.

If, additionally, $\det \widetilde{M} \neq 0$, then we use the contradiction argument to show that the condition (\ref{eq:Finite_eq_15}) is necessary for (\ref{eq:Finite_eq_1}) to be memory-type null controllable. Assume that (\ref{eq:Finite_eq_1}) is memory-type null controllable but the condition (\ref{eq:Finite_eq_15}) does not hold. Then, in view of (\ref{eq:Finite_eq_26}),

\[
\text{rank} \begin{bmatrix}
B & A_1 B & A_2 B & \cdots & A_i B & A_{i+1} B & \cdots \\
0 & B & A_1 B & \cdots & A_{i-1} B & A_i B & \cdots
\end{bmatrix} < 2n.
\]


This implies that there is a $(w_T, z_T) \in \mathbb{R}^{2n} \setminus \{0\}$ satisfying
\begin{equation}\label{eq:Finite_eq_32}
(w_T^\top, -z_T^\top \widetilde{M}) 
\begin{pmatrix}
B & A_1 B & A_2 B & \cdots & A_i B & A_{i+1} B & \cdots \\
0 & B & A_1 B & \cdots & A_{i-1} B & A_i B & \cdots
\end{pmatrix} = 0.
\end{equation}

Clearly, this $(w_T, z_T)$ satisfies
\[
\left\{
\begin{aligned}
& B^\top w_T = 0, \\
& B^\top A_1^\top w_T - B^\top \widetilde{M}^\top z_T = 0, \\
& B^\top A_i^\top w_T - B^\top A_{i-1}^\top \widetilde{M}^\top z_T = 0,\quad i = 2, 3, \ldots
\end{aligned}
\right.
\]

Hence, the corresponding solution $w(\cdot)$ of (\ref{eq:Finite_eq_2}) satisfies
\begin{equation}\label{eq:Finite_eq_33}
\left. \frac{d^k B^\top w(t)}{dt^k} \right|_{t=T} = 0, \quad k = 1, 2, \ldots
\end{equation}

Since $B^\top w(\cdot)$ is an analytic function, (\ref{eq:Finite_eq_33}) implies that $B^\top w(\cdot) \equiv 0$. In view of Theorem 3.1, this leads to $w_T = z_T = 0$, a contradiction. 






    
\end{proof}











XXXXXXXXXXXXXXXXXXXXXXXXXXXXXXXXXXXXXXXXXXXXXX
\begin{proposition}
\label{prop:2.1}
The system described by \eqref{eq:Intro_eq_7} is Delay and memory-type null controllable (with the kernel $\widetilde{M}(\cdot)$) if and only if every solution to \eqref{eq:dual_eq_1} satisfies the estimate
\begin{equation}
\label{eq:dual_eq_2}
|w(\theta)|^2_Y \leq K \int_0^T \|B(t)^*w(t)\|^2_U\,dt \quad \forall \, w_T, z_T \in Y.
\end{equation}
\end{proposition}

\begin{proof}
The proof follows standard arguments, and we include it here for completeness.

To prove the "if" direction, let us fix any $\phi(\theta) \in Y$. Define a linear subspace $\mathcal{L} \subset L^2(0, T; U)$ by
\[
\mathcal{L} = \left\{ B(\cdot)^*w(\cdot) \,\middle|\, w(\cdot) \text{ is a solution to } \eqref{eq:dual_eq_1} \text{ for some } w_T, z_T \in Y \right\}.
\]

For each $B(\cdot)^*w(\cdot) \in \mathcal{L}$, we define the linear functional
\[
\widetilde{\mathcal{L}}(B(\cdot)^*w(\cdot)) = -\langle w(\theta), \phi(\theta) \rangle_Y.
\]

Using inequality \eqref{eq:dual_eq_2}, it follows that $\widetilde{\mathcal{L}}$ is a bounded linear functional on the normed space $\mathcal{L}$, where the norm is inherited from $L^2(0, T; U)$. By the Hahn–Banach theorem, $\widetilde{\mathcal{L}}$ can be extended to a bounded linear functional on all of $L^2(0, T; U)$. Then, by the Riesz representation theorem, there exists a function $\eta(\cdot) \in L^2(0, T; U)$ such that

\begin{equation}
\label{eq:dual_eq_3}
\int_0^T \langle B(t)^*w(t), \eta(t) \rangle_U \, dt = -\langle w(\theta), \phi(\theta) \rangle_Y.
\end{equation}

We now assert that the control we seek is given by
\begin{equation}
\label{eq:dual_eq_4}
u(\cdot) = \eta(\cdot).
\end{equation}

To verify this, consider any $w_T, z_T \in Y$. Using \eqref{eq:Intro_eq_7} and \eqref{eq:dual_eq_1}, we derive:
\begin{align}
\label{eq:dual_eq_5}
&\langle y(T), w_T \rangle_Y - \langle w(\theta), \phi(\theta) \rangle_Y \notag \\
&= \int_0^T \left\langle \frac{d}{dt} y(t), w(t) \right\rangle_Y \, dt + \int_0^T \left\langle y(t), \frac{d}{dt} w(t) \right\rangle_Y \, dt \notag \\
&= \int_0^T \left\langle Ay(t) + A_1y(t-h) + \int_0^t M(t-s)y(s)\,ds + B(t)u(t), w(t) \right\rangle_Y \, dt \notag \\
&\quad + \int_0^T \left\langle y(t), -A^*w(t) - A_1w(t-h) - \int_t^T M(s - t)^*w(s)\,ds + \widetilde{M}(T - t)^*z_T \right\rangle_Y \, dt \notag \\
&= \int_0^T \langle B(t)^*w(t), u(t) \rangle_U \, dt + \int_0^T \left\langle y(t), \widetilde{M}(T - t)^*z_T \right\rangle_Y \, dt.
\end{align}
\footnotetext{Throughout this paper, for any operator-valued function $R$, we denote its pointwise adjoint by $R^*$. For instance, if $R \in L^1(0,T; \mathcal{L}(Y))$, then $R^* \in L^1(0,T; \mathcal{L}(Y))$ and satisfies $\|R\|_{\mathcal{L}(Y)} = \|R^*\|_{\mathcal{L}(Y)}$.}

By combining \eqref{eq:dual_eq_3}, \eqref{eq:dual_eq_4}, and \eqref{eq:dual_eq_5}, we obtain
\begin{equation}
\left( w_T, y(T) \right)_Y - \left( z_T, \int_0^T \widetilde{M}(T - t) y(t)\,dt \right)_Y = 0 \quad \forall w_T, z_T \in Y.
\end{equation}
Consequently, $y(T) = \int_0^T \widetilde{M}(T - t) y(t)\,dt = 0$, which confirms the desired result.

We now proceed to prove the “only if” direction. Given any $w_T, z_T \in Y$, and using \eqref{eq:dual_eq_1}, we define a bounded linear operator $\mathcal{F}: Y \times Y \to Y$ as follows:
\begin{equation}
\label{eq:dual_eq_6}
\mathcal{F}(w_T, z_T) := w(\theta),
\end{equation}
where $w(t)$ denotes the solution to the adjoint system \eqref{eq:dual_eq_1}.

To prove \eqref{eq:dual_eq_2}, we use a contradiction argument. Suppose \eqref{eq:dual_eq_2} is false. Then there exist sequences $\{w_T^k\}, \{z_T^k\} \subset Y$ such that the corresponding solutions $w^k(\cdot)$ to \eqref{eq:dual_eq_1} (with $(w_T, z_T)$ replaced by $(w_T^k, z_T^k)$) satisfy:
\begin{equation}
\label{eq:dual_eq_7}
0 \leq \int_0^T \|B(t)^*w^k(t)\|_U^2\,dt \leq \frac{1}{k^2} \|w^k(\theta)\|_Y^2, \quad \forall k \in \mathbb{N}.
\end{equation}

Define the normalized vectors as follows:
\[
\widehat{w}_T^k := \frac{w_T^k}{\|w^k(\theta)\|_Y}, \quad \widehat{z}_T^k := \frac{z_T^k}{\|w^k(\theta)\|_Y},
\]
and let $\widehat{w}^k(\cdot)$ denote the corresponding solution to \eqref{eq:dual_eq_1} for $(w_T, z_T) = (\widehat{w}_T^k, \widehat{z}_T^k)$. Then, using \eqref{eq:dual_eq_6} and \eqref{eq:dual_eq_7}, we obtain:
\begin{equation}
\label{eq:dual_eq_8}
\int_0^T \|B(t)^*\widehat{w}^k(t)\|_U^2\,dt \leq \frac{1}{k^2}, \quad \|\widehat{w}^k(\theta)\|_Y = 1.
\end{equation}

Since \eqref{eq:Intro_eq_7} is assumed to be Delay and memory-type null controllable (with memory kernel $\widetilde{M}(\cdot)$), for any $y_0 \in Y$, there exists a control $u(\cdot) \in L^2(0,T;U)$ such that the associated solution $y(\cdot)$ satisfies \eqref{eq:Intro_eq_9}. Then, for any $w_T, z_T \in Y$, a derivation similar to \eqref{eq:dual_eq_5} gives:
\[
-(w(\theta), \phi(\theta))_Y = \int_0^T (B(t)^*w(t), u(t))_U\,dt.
\]

In particular, for the normalized sequences, this leads to:
\begin{equation}
\label{eq:dual_eq_9}
-\left( \mathcal{F}(\widehat{w}_T^k, \widehat{z}_T^k), y_0 \right)_Y = \int_0^T (B(t)^*\widehat{w}^k(t), u(t))_U\,dt.
\end{equation}

From \eqref{eq:dual_eq_9} and the first part of \eqref{eq:dual_eq_8}, it follows that $\mathcal{F}(\widehat{w}_T^k, \widehat{z}_T^k)$ converges weakly to zero in $Y$. Then, by the Banach-Steinhaus theorem (principle of uniform boundedness), the sequence $\{ \mathcal{F}(\widehat{w}_T^k, \widehat{z}_T^k) \}$ is uniformly bounded in $Y$. However, this contradicts the normalization $\|\widehat{w}^k(\theta)\|_Y = 1$ in \eqref{eq:dual_eq_8}.

This contradiction completes the proof of Proposition~\ref{prop:2.1}.
\end{proof}
XXXXXXXXXXXXXXXXXXXXXXXXXXXXXXXXXXX\\


